\Needspace{10\baselineskip}
\SourceLeaf{338}

\begin{center}
[\emph{Lecture on the Theory of Imaginary Quantities.}]

On the Quadrilateral.

[1.]
\end{center}

Given a quadrilateral (1), (2), (3), (4) determined by the complex numbers $p'$, $p''$, $p'''$,
$p^{\mathrm{IV}}$, a general formula can be obtained by means of the laterals
by passing to a triangle (I), (II), (III) derived from it in the only
symmetric way, its vertices being fixed as follows.
\[
\begin{array}{rll}
\text{(I)} & \text{by} & p'p'' + p'''p^{\mathrm{IV}} \\
\text{(II)} & \text{,,} & p'p''' + p'' p^{\mathrm{IV}} \\
\text{(III)} & \text{,,} & p'p^{\mathrm{IV}} + p'' p'''.
\end{array}
\]
Suppose, for\,example, that the vertices are fixed by the following complex numbers.
\[
\begin{array}{rlll}
\text{(1)} & \text{is fixed by} & : & 2 + \sqrt{-1} = p', \\
\text{point (2)} & \text{,,} & : & 3 \hspace{2.8em} = p'', \\
\text{,, (3)} & \text{,,} & : & 1 + 2\sqrt{-1} = p''', \\
\text{,, (4)} & \text{,,} & : & 1 + 3\sqrt{-1} = p^{\mathrm{IV}},
\end{array}
\]
Then
\[
\begin{array}{rlll}
\text{(I)} & \text{is determined by} & : & 1 + 8\sqrt{-1} \\
\text{(II)} & \text{,, ,,} & : & 3 + 14\sqrt{-1} \\
\text{(III)} & \text{,, ,,} & : & 2 + 13\sqrt{-1}.
\end{array}
\]

Referring (II) to (I), I obtain, by subtracting the complex numbers corresponding to these vertices
from one another,
\[
\text{(II)} - \text{(I)} = p'(p''' - p'') + p^{\mathrm{IV}}(p'' - p'''),
\]
and similarly
\[
\text{(III)} - \text{(II)} = p'(p^{\mathrm{IV}} - p''') + p'' (p''' - p^{\mathrm{IV}})
\]
\[
\text{(I)} - \text{(III)} = p'(p'' - p^{\mathrm{IV}}) + p'''(p^{\mathrm{IV}} - p'')
\]
or
\[
\begin{array}{rl}
\text{(II)} - \text{(I)} & = (p' - p^{\mathrm{IV}})(p''' - p'') \\
\text{(III)} - \text{(II)} & = (p' - p'')(p^{\mathrm{IV}} - p''') \\
\text{(I)} - \text{(III)} & = (p' - p''')(p'' - p^{\mathrm{IV}}).
\end{array}
\]

Referring the angles of the quadrilateral to a fixed axis and denoting, for\,example, the angle that (4,1)
makes with it by [4,1], I obtain
\[
\begin{array}{rl}
-p^{\mathrm{IV}} + p' & = (4,1)\{\cos[4,1] + i\sin[4,1]\} \\
p''' - p'' & = (3,2)\{\cos[3,2] + i\sin[3,2]\}
\end{array}
\]
\[
\text{etc.}
\]
Hence we have
\[
\begin{array}{rl}
\text{(I, II)} & = (4,1) . (3,2) \\
\text{(II, III)} & = (2,1) . (4,3) \\
\text{(III, I)} & = (3,1) . (4,2).
\end{array}
\]
